\section{Función descriptiva fraccionaria de Machado y transferencia sesgada}
\label{sec:machado-fdf}

\HAFOCardMachado

Esta sección responde dos preguntas que deben mantenerse separadas. La primera
es cómo entra un desplazamiento continuo en una realización de Lur'e; la
segunda es qué significa elevar una función descriptiva a un orden real en el
sentido propuesto por Machado \cite{Machado2015FDF}. Ninguna de las dos
operaciones cambia por sí sola el orden de Caputo del bloque dinámico.

\subsection{Transferencia, sesgo y transferencia incremental}
\label{sec:biased-transfer-distinction}

Considérese la realización conmensurable
\begin{equation}
 {}^{\mathrm C}D_{0+}^{q}x(t)
 =Px(t)+b\,\psi(\sigma(t)),
 \qquad
 \sigma(t)=r^{\mathsf T}x(t),
 \qquad 0<q\leq1.
 \label{eq:machado-lure}
\end{equation}
La transferencia entrada--salida del bloque lineal es
\begin{equation}
 W_q(s)=r^{\mathsf T}(s^qI-P)^{-1}b=G(s^q).
 \label{eq:machado-Wq}
\end{equation}
Para el predictor armónico se fija la rama principal
\begin{equation}
 (\mathrm i\omega)^q
 =\omega^q
 \left[
 \cos\!\left(\frac{\pi q}{2}\right)
 +\mathrm i\sin\!\left(\frac{\pi q}{2}\right)
 \right],
 \qquad \omega>0.
 \label{eq:machado-iomegaq}
\end{equation}

Un sesgo en la entrada no crea automáticamente otra transferencia. Si
\begin{equation}
 \sigma(t)=\sigma_0+A\cos(\omega t),
 \label{eq:biased-input}
\end{equation}
el bloque lineal continúa siendo \(W_q\); lo que cambia son el valor medio y
el fundamental de \(\psi(\sigma_0+A\cos\vartheta)\). Por eso, en el resto del
informe, la expresión ``transferencia sesgada'' debe leerse como abreviatura
de \emph{cierre con función descriptiva sesgada}, no como un nuevo operador
lineal.

Existe, sin embargo, un tercer objeto legítimo. Si \(\psi\) es diferenciable
en un punto de operación \(\bar\sigma\), la linealización incremental es
\begin{equation}
 P_{\bar\sigma}=P+b\psi'(\bar\sigma)r^{\mathsf T},
 \qquad
 W_{q,\bar\sigma}^{\mathrm{inc}}(s)
 =r^{\mathsf T}(s^qI-P_{\bar\sigma})^{-1}b.
 \label{eq:incremental-biased-transfer}
\end{equation}
Esta transferencia describe perturbaciones pequeñas alrededor del punto de
operación. No debe usarse simultáneamente con una función descriptiva que ya
contenga la misma pendiente, pues se contaría dos veces. En una característica
PWA sobre una frontera se necesita especificar pendiente lateral, derivada de
Clarke o convención de Filippov; sin esa declaración
\cref{eq:incremental-biased-transfer} no es un objeto único.

\subsection{Dos fraccionalidades distintas}

La fraccionalidad de \cref{eq:machado-Wq} pertenece al bloque dinámico: se
sustituye \(s\) por \(s^q\). Es distinta de la
\emph{función descriptiva fraccionaria} propuesta por Machado. En nuestra
notación, la operación central de esa propuesta es
\begin{equation}
 {}^{\mu}\!N(A,\omega)
 =\bigl[N(A,\omega)\bigr]^\mu.
 \label{eq:machado-published}
\end{equation}
El artículo original estudia la potencia de funciones descriptivas complejas,
principalmente para backlash, propone una interpretación análoga a elementos
fraccionarios y discute una realización aproximada mediante asociaciones de
elementos con backlash \cite{Machado2015FDF}. La realizabilidad general no se
establece como teorema. Por tanto, \cref{eq:machado-published} no transforma
automáticamente una no linealidad estática cualquiera en una no linealidad
fraccionaria físicamente equivalente.

En particular,
\begin{equation}
 q\neq\mu\quad\text{en general}.
 \label{eq:q-not-mu}
\end{equation}
El parámetro \(q\) es el orden de la derivada del sistema; \(\mu\) recorre una
familia de ganancias armónicas. No hay fundamento para imponer \(\mu=q\), ni
aparece en el cierre un producto de órdenes \(q\mu\).

\subsection{Función descriptiva física y convención fasorial}

Para una entrada centrada
\[
 \sigma(t)=A\cos(\omega t),\qquad A>0,
\]
sea \(Y_1(A,\omega)\) el fasor fundamental de la salida. Adoptamos
\begin{equation}
 Y_1(A,\omega)
 =\frac{1}{\pi}\int_{0}^{2\pi}
 y(\vartheta;A,\omega)e^{-\mathrm i\vartheta}\,d\vartheta,
 \qquad
 N(A,\omega)=\frac{Y_1(A,\omega)}{A}.
 \label{eq:machado-classical-df}
\end{equation}
Si \(\psi\) es estática, impar y sin memoria,
\(y(\vartheta;A)=\psi(A\cos\vartheta)\), la parte seno se anula y
\begin{equation}
 N(A)=\frac{1}{\pi A}\int_{0}^{2\pi}
 \psi(A\cos\vartheta)\cos\vartheta\,d\vartheta\in\mathbb R.
 \label{eq:machado-real-df}
\end{equation}
Si hay histéresis, impacto o un estado interno, \(Y_1\) puede depender de
\(\omega\), de la órbita periódica interna y de su prehistoria. Sustituirlo por
\cref{eq:machado-real-df} sería entonces una pérdida de modelo.

\subsection{Potencias complejas, hojas y unidades}

Sea
\[
 N(A,\omega)=\rho(A,\omega)e^{\mathrm i\phi(A,\omega)},
 \qquad \rho>0.
\]
La potencia compleja es multivaluada:
\begin{equation}
 \bigl[N(A,\omega)\bigr]_{\ell}^{\mu}
 =\rho^\mu
 \exp\!\left\{\mathrm i\mu[\phi+2\pi\ell]\right\},
 \qquad \ell\in\mathbb Z.
 \label{eq:machado-multivalued}
\end{equation}
La rama principal es
\begin{equation}
 \bigl[N(A,\omega)\bigr]_{\mathrm{pr}}^{\mu}
 =\exp\!\left[\mu\operatorname{Log}N(A,\omega)\right],
 \qquad \operatorname{Arg}N\in(-\pi,\pi].
 \label{eq:machado-principal}
\end{equation}
Una continuación no debe aplicar la rama principal punto a punto cuando la
curva cruza \(( -\infty,0]\). Deben desenvolverse la fase, fijarse una hoja y
registrarse los cruces del corte.

Para un DF real hay tres casos:
\begin{enumerate}
 \item Si \(N>0\), existe una rama real positiva única.
 \item Si \(N=0\) y \(\mu>0\), puede extenderse por continuidad con
       \(0^\mu=0\), pero \(\operatorname{Log}N\) no existe y la sensibilidad
       es singular para \(0<\mu<1\).
 \item Si \(N<0\) y \(\mu\notin\mathbb Z\), la rama principal es compleja.
       Algunos exponentes racionales con denominador impar admiten una raíz
       real, que no tiene por qué coincidir con la rama principal.
\end{enumerate}

También existe una restricción dimensional. Si \(N\) tiene unidades,
\(N^\mu\) cambia las unidades. Con una escala \(N_\star\) de las mismas
unidades se define
\begin{equation}
 N_{\mu,N_\star}(A,\omega)
 =N_\star\left[\frac{N(A,\omega)}{N_\star}\right]^\mu.
 \label{eq:machado-normalized}
\end{equation}
La normalización adoptada en los ejemplos equivale a \(N_\star=1\); presupone,
por tanto, que la función descriptiva fue adimensionalizada.

\subsection{Familia auxiliar normalizada sobre la rama positiva}

Sobre la rama real positiva se define
\begin{equation}
 N_\mu^{(+)}(A)=N(A)^\mu,
 \qquad N(A)>0,\quad \mu>0,
 \label{eq:hafo-machado}
\end{equation}
para \(N(A)>0\) y \(\mu>0\). Esta familia es aplicable al Chua no suave bajo
esa restricción; quedan fuera \(N(A)\leq0\), \(\mu\leq0\), una función
descriptiva compleja y la rama arctangente aquí considerada. Para
\(\mu\neq1\),
\begin{equation}
 N_\mu^{(+)}(A)
 \neq\frac{1}{\pi A}\int_0^{2\pi}
 \psi(A\cos\vartheta)\cos\vartheta\,d\vartheta
 \label{eq:machado-not-physical-fourier}
\end{equation}
en general. En consecuencia, el nombre preciso es
\emph{familia auxiliar inspirada en Machado}: modifica la ecuación de
amplitud y diversifica semillas, pero no es el coeficiente de Fourier físico
de la no linealidad original.

\subsection{Cierre centrado y dos residuos obligatorios}

El cierre físico de primer armónico es
\begin{equation}
 R_{\mathrm{fis}}(A,\omega;q)
 =1-W_q(\mathrm i\omega)N(A,\omega)=0.
 \label{eq:machado-physical-closure}
\end{equation}
La familia auxiliar usa
\begin{equation}
 R_{\mu}(A,\omega;q)
 =1-W_q(\mathrm i\omega)N_{\mu,N_\star}(A,\omega)=0.
 \label{eq:machado-aux-closure}
\end{equation}
Ambos coinciden para \(\mu=1\), no para \(\mu\neq1\). Si \(N(A)>0\) es
real, el cierre auxiliar se separa en
\begin{align}
 \operatorname{Im}W_q(\mathrm i\omega)&=0,
 \label{eq:machado-frequency}\\
 k(\omega)=\frac{1}{\operatorname{Re}W_q(\mathrm i\omega)}&>0,
 \label{eq:machado-gain}\\
 N(A)&=N_\star\left(\frac{k}{N_\star}\right)^{1/\mu}.
 \label{eq:machado-amplitude}
\end{align}
Para \(N_\star=1\), se resuelve \(N(A)=k^{1/\mu}\). Además de verificar que
el objetivo pertenece al rango de \(N\), se debe controlar
\begin{equation}
 \frac{dA}{d\mu}
 =-\frac{N(A)\log[N(A)/N_\star]}{\mu N'(A)},
 \label{eq:machado-sensitivity}
\end{equation}
pues la raíz está mal condicionada cerca de \(N'(A)=0\), \(N=0\) o un
extremo de la curva de ganancia.

Cada candidato debe conservar, sin mezclarlos,
\begin{align}
 \rho_{\mathrm{aux}}
 &=\left|1-W_q(\mathrm i\omega)N_{\mu,N_\star}(A,\omega)\right|,
 \label{eq:machado-rho-aux}\\
 \rho_{\mathrm{fis}}
 &=\left|1-W_q(\mathrm i\omega)N(A,\omega)\right|.
 \label{eq:machado-rho-phys}
\end{align}
Un \(\rho_{\mathrm{aux}}\) pequeño no implica un \(\rho_{\mathrm{fis}}\)
pequeño.

Para reconstruir una semilla sustituta se forma
\begin{equation}
 P_k=P+bkr^{\mathsf T},\qquad
 \lambda=(\mathrm i\omega)^q,\qquad
 P_kv=\lambda v,\qquad r^{\mathsf T}v=1,
 \label{eq:machado-modal}
\end{equation}
y
\begin{equation}
 x_{\mathrm{seed}}=A\operatorname{Re}(ve^{\mathrm i\theta_0}).
 \label{eq:machado-seed}
\end{equation}
Es una semilla del sustituto linealizado por \(k\), no una solución armónica
exacta del PVI causal de Caputo.

\subsection{Cierre sesgado de Machado}

Para \cref{eq:biased-input}, los coeficientes físicos son
\begin{align}
 \Psi_0(\sigma_0,A)
 &=\frac{1}{2\pi}\int_0^{2\pi}
 \psi(\sigma_0+A\cos\vartheta)\,d\vartheta,
 \label{eq:machado-biased-dc}\\
 Y_1(\sigma_0,A)
 &=\frac{1}{\pi}\int_0^{2\pi}
 \psi(\sigma_0+A\cos\vartheta)e^{-\mathrm i\vartheta}\,d\vartheta,
 \label{eq:machado-biased-y1}\\
 N_b(\sigma_0,A)&=\frac{Y_1(\sigma_0,A)}{A}.
 \label{eq:machado-biased-df}
\end{align}
El cierre continuo permanece físico:
\begin{equation}
 \sigma_0-W_q(0)\Psi_0(\sigma_0,A)=0.
 \label{eq:machado-biased-dc-closure}
\end{equation}
Una familia auxiliar sesgada posible es
\begin{equation}
 N_{\mu,b,N_\star}(\sigma_0,A)
 =N_\star\left[\frac{N_b(\sigma_0,A)}{N_\star}\right]^\mu,
 \label{eq:machado-biased-family}
\end{equation}
con hoja registrada, y el cierre sustituto
\begin{equation}
 1-W_q(\mathrm i\omega)N_{\mu,b,N_\star}(\sigma_0,A)=0.
 \label{eq:machado-biased-closure}
\end{equation}
No se eleva \(\Psi_0\) a \(\mu\): es el valor medio físico que entra en el
equilibrio afín. También se evalúa
\begin{equation}
 \rho_{\mathrm{fis},b}
 =\left|1-W_q(\mathrm i\omega)N_b(\sigma_0,A)\right|.
 \label{eq:machado-biased-physical-residual}
\end{equation}
Las evaluaciones sesgadas validadas en este volumen usan \(\mu=1\); coinciden
con la familia sesgada clásica y no validan todavía
\cref{eq:machado-biased-family} para \(\mu\neq1\).

\subsection[Ejemplo Wolfram: Chua PWL, q=0.9998]
{Ejemplo Wolfram: Chua PWL, \(q=0.9998\)}

Para
\[
 \alpha=8.4562,\quad \beta=12.0732,\quad \gamma=0.0052,
 \quad m_1=-1.1468,
 \quad \Delta=m_0-m_1=0.97,
\]
se usa
\begin{equation}
 P=\begin{pmatrix}
 -\alpha(1+m_1)&\alpha&0\\
 1&-1&1\\
 0&-\beta&-\gamma
 \end{pmatrix},\qquad
 b=\begin{pmatrix}-\alpha\\0\\0\end{pmatrix},\qquad
 r=\begin{pmatrix}1\\0\\0\end{pmatrix}.
 \label{eq:machado-chua-realization}
\end{equation}
La no linealidad residual es
\(\psi(\sigma)=\Delta\operatorname{sat}(\sigma)\), con
\begin{equation}
 N(A)=\begin{cases}
 \Delta,&0<A\leq1,\\[1mm]
 \displaystyle\frac{2\Delta}{\pi}
 \left[\arcsin\!\left(\frac1A\right)
 +\frac{\sqrt{A^2-1}}{A^2}\right],&A>1.
 \end{cases}
 \label{eq:machado-chua-df}
\end{equation}
Wolfram confirma
\begin{equation}
 N'(A)=-\frac{4\Delta\sqrt{A^2-1}}{\pi A^3}<0,
 \qquad A>1.
 \label{eq:machado-chua-monotone}
\end{equation}

En la rama \(\omega_0=2.0402860510795615\), la evaluación independiente
produjo
\begin{align}
 W_q(\mathrm i\omega_0)
 &=4.761387970782572-1.52\times10^{-13}\mathrm i,\\
 k&=0.21002279296212067.
\end{align}
El carril clásico \(\mu=1\) da
\(A_{\mathrm{cl}}=5.851767785486093\). Para \(\mu=2\), HAFO registra
\begin{equation}
 A_2=2.628433514775448,
 \quad N(A_2)=0.458282437981291,
 \quad N(A_2)^2=0.21002279296207582.
\end{equation}
Por tanto,
\begin{align}
 \rho_{\mathrm{aux}}&=2.16\times10^{-13},\\
 \rho_{\mathrm{fis}}&=1.182060487425029.
 \label{eq:machado-two-residuals-example}
\end{align}
El cierre sustituto es preciso, pero el balance físico del Chua original no
se satisface. La rama \(\mu=2\) sólo justifica una condición inicial diferente
para la integración causal posterior.

\begin{warningbox}
Un único ``residuo armónico'' que incluya también errores de reconstrucción
modal no sustituye los dos residuos algebraicos de
\cref{eq:machado-rho-aux,eq:machado-rho-phys}. Los tres valores deben
reportarse por separado y con nombres inequívocos.
\end{warningbox}

\subsection{Casos convenientes, casos riesgosos y protocolo}

La familia de Machado es especialmente pertinente cuando:
\begin{enumerate}
 \item existe backlash, fricción o histéresis con DF complejo modelado y fase
       continua \cite{DuarteMachado2009Coulomb,Machado2015FDF};
 \item la ganancia SISO es adimensional o está normalizada;
 \item \(N\) es real, positiva y monótona en el intervalo explorado;
 \item se desea diversificar amplitudes antes de la continuación causal;
 \item la malla en \(\mu\) está anclada en \(\mu=1\) y recibe el mismo
       presupuesto que las semillas de control.
\end{enumerate}
Debe evitarse o requiere teoría adicional si \(N\) cambia de signo, cruza el
corte de rama, tiene unidades sin \(N_\star\), la realización es MIMO, hay
resonancia armónica intensa, impactos o una respuesta de banda ancha. En una
matriz de funciones descriptivas, una potencia funcional matricial no es una
potencia elemento a elemento; exige espectro, ramas y posible tratamiento de
bloques de Jordan.

Para Chua PWL centrado, \cref{eq:hafo-machado} es admisible como diversificador
de semillas. Aplicarla al Chua arctangente, Kalman--Fitts o MAVPD sería una
extensión que requeriría validación adicional, no una consecuencia del artículo
de Machado. En PLL, \(2J_1(A)/A\) cambia de signo y tiene ceros: sólo pueden
usarse intervalos positivos separados o una rama compleja explícita.

El protocolo completo es:
\begin{enumerate}
 \item derivar \(N\), \(\Psi_0\) y \(N_b\) desde integrales de Fourier;
 \item declarar \(q,\mu,N_\star\), hoja logarítmica y desenvolvimiento de fase;
 \item resolver \(\mu=1\) como control;
 \item guardar \(\rho_{\mathrm{aux}}\) y \(\rho_{\mathrm{fis}}\);
 \item medir distorsión armónica física, por ejemplo
 \[
 H_{\mathrm{NL}}=
 \frac{(\sum_{n\ge2}|Y_n|^2)^{1/2}}{|Y_1|},
 \]
 y su versión ponderada por \(|W_q(\mathrm in\omega)|\);
 \item descartar raíces mal condicionadas y cruces de rama no controlados;
 \item integrar el sistema original, no el sustituto, con igual terminal
       inferior, prehistoria, paso, horizonte y política de memoria;
 \item refinar paso y horizonte y comparar con semillas clásica, sesgada,
       multiharmónica, perpetua y aleatoria;
 \item ejecutar después clasificación dinámica y sondeos de cuenca.
\end{enumerate}

El balance con \(s^q\) sigue siendo un predictor estacionario de tipo
Liouville--Weyl. En un PVI causal de Caputo,
\[
 \mathcal L\{{}^{\mathrm C}D^qx\}
 =s^qX(s)-s^{q-1}x(0),
\]
de modo que una semilla sinusoidal no es por sí misma una solución periódica
exacta desde \(t=0\). Machado complementa la función descriptiva clásica, el
cierre sesgado y la continuación; no reemplaza la memoria causal ni las
pruebas de cuenca.
